\chapter{Quantum Gravity Candidates}
\label{chap:quantum-gravity-candidates}

Since the ultimate 'Theory of Everything' remains far out of reach and difficult to even define, we may have to settle for a more modest milestone: unifying General Relativity and Quantum Mechanics.

\section{Semiclassical Gravity}

One of the earliest attempts to bridge the quantum--classical divide is \emph{semiclassical gravity}. In this approach, matter is treated as fully quantum, while spacetime remains classical. To make the Einstein field equations workable, the operator-valued stress--energy tensor of quantum matter is replaced by its expectation value—the renormalized average of the energy and momentum calculated over the quantum state of the matter fields. This resulting set of ordinary numbers can then be inserted into the equations governing curvature.

Semiclassical gravity is remarkably successful. It accurately describes a wide range of phenomena, from laboratory experiments to astrophysical observations and cosmology. It even predicts striking effects such as Hawking radiation in black holes. Yet its very success also exposes its conceptual limitation: the approach is \emph{ad hoc}.
The theory works well for everything we can observe, but it does not answer any of the deeper question, like what is the physics at the singularity of a black hole.

\section{Perturbative Quantum Gravity}

A natural next step is \emph{perturbative quantum gravity}, where spacetime is expanded around a simple background—typically flat or slightly curved—and the perturbations are treated as quantum fields. This approach is conceptually straightforward and extends the familiar machinery of quantum field theory to gravity.

However, it quickly runs into a fundamental problem: gravity is nonrenormalizable. Unlike the Standard Model, where quantum infinities can be systematically controlled, perturbative gravity's divergences become increasingly severe at higher energies. To cancel these, we are forced to introduce new counterterms at every order of the calculation. This process never terminates, requiring an infinite number of independent parameters to be measured and fixed. Ultimately, this strips the theory of its predictive power; instead of a concise model, we are left with an infinite list of unknowns. The techniques that work spectacularly well for matter fields simply break down for spacetime itself.

\section{Nonperturbative and Geometric Approaches}

In response to the failure of perturbative quantization, researchers have developed nonperturbative frameworks that do not assume a fixed background geometry. These include Loop Quantum Gravity (LQG), Causal Dynamical Triangulations, Asymptotic Safety, Causal Set Theory and Noncommutative Geometry (NCG). While all these represent valuable research programs, they all remain incomplete as full theories of quantum gravity.


\subsection*{Loop Quantum Gravity}

A leading example is Loop Quantum Gravity (LQG), which models spacetime as a discrete combinatorial structure of spin networks. LQG is mathematically rigorous and fully background-independent, offering a conceptually clean quantization of geometry.

The theory is a cousin to Penrose’s Twistor program in that his work on spinors and spin networks provided the mathematical toolbox that allowed LQG to exist.

However, it seems major obstacles still remain. Deriving a smooth classical spacetime limit is nontrivial, and embedding standard particle physics into the LQG framework remains unresolved.
Despite its mathematical rigor and background independence, LQG faces persistent challenges in recovering a realistic low-energy limit. The emergence of a smooth, 4-dimensional classical spacetime with the correct long-distance dynamics (including the Einstein equations) is not fully under control, often relying on semiclassical approximations or specific choices of states. Incorporating matter fields, particularly chiral fermions from the Standard Model, encounters issues like fermion doubling or difficulties in coupling them consistently without breaking key symmetries. Lorentz invariance violations predicted in some early formulations have not been observed experimentally, forcing adjustments that weaken distinctive predictions. Overall, while it quantizes geometry cleanly, it has not yet produced a complete, predictive quantum theory of gravity unified with particle physics.




\section{Canonical Quantum Gravity}

A conceptually direct attempt to quantize gravity involves applying canonical quantization to General Relativity. In this approach, 4D spacetime is decomposed into a foliation of spatial slices (the 3+1 formalism), allowing the Einstein equations to be rewritten in Hamiltonian form, for compatibility with Quantum Mechanics.

Upon quantization, the classical constraints of General Relativity are promoted to operator equations acting on a wavefunction, which encodes the quantum state of the spatial geometry.
The central result is the Wheeler-DeWitt equation:

\[
\hat{H}\Psi[h_{ij}] = 0.
\]

The irony is that while trying to make General Relativity compatible with Schrödinger equation, the resulting equation actually lacks the time derivative essential in the Schrödinger equation. By trying to force GR into the "Hamiltonian" box of QM, the time variable disappears entirely. 

This leads to the "Problem of Time": whereas standard quantum mechanics describes states evolving against a background time, the wavefunction of the universe appears static. Recovering time as an emergent or relational concept remains a profound open problem.

Stephen Hawking’s Euclidean Quantum Gravity program utilized the Wheeler-DeWitt framework alongside Wick rotations and path integrals to define the "No-Boundary" proposal.


\section{String Theory: A Rock Too Heavy?}

In string theory, the traditional concept of zero-dimensional point particles is replaced by one-dimensional strings.
The vibrational frequency of a string determines the mass and charge of a particle, much like different notes played on a guitar string.

The theory's primary selling point is its inherent ability to unify gravity with the standard model of particle physics.
Gravity emerges naturally as one of the vibrational modes of the string—specifically the massless, spin-2 graviton.
Beyond this, the framework encompasses deep mathematical structures such as dualities, extra dimensions, and black-hole entropy counting.

The earliest framework, bosonic string theory, is mathematically consistent only in 26 spacetime dimensions.
Later, the introduction of supersymmetry led to superstring theories, which include fermionic degrees of freedom and are consistent in 10 dimensions.
These dimensions are not arbitrary; consistency requires the cancellation of quantum anomalies on the string worldsheet.
From a mathematical perspective, these constraints represent a form of elegance: the structure of the theory tightly restricts the space of consistent possibilities.

However, several open issues prevent string theory from being accepted as a complete physical theory.

First, most formulations are not manifestly background independent, because it doesn't explain but pre-assumes spacetime geometry.
If a theory starts with a stage already built, someone is entitled to ask who hired the carpenter.
This stands in contrast to general relativity!

Second, the theory relies on supersymmetry (SUSY), which predicts that every known particle has a superpartner with a spin differing by $1/2$.
Despite the high energy reach of modern colliders, none of these predicted sparticles have been detected.

Third, the theory admits an enormous landscape of possible vacuum states, estimated at $10^{500}$. To put this in perspective, the number of atoms (baryons) in the universe
is estimated to be $10^{80}$.

Our universe would correspond to the one that happens to produce the physics we observe.
These vacua arise from the myriad ways of compactifying the extra dimensions into complex shapes known as Calabi-Yau manifolds.
The specific geometry of these manifolds, combined with the way higher-dimensional objects called ``branes'' wrap around them, dictates the physics of the resulting 4D universe.
This raises significant concerns regarding predictivity and falsifiability; if the theory can accommodate almost any physics, it may lack the power to uniquely predict our own.

Despite these challenges, string theory remains a dominant tool for theoretical discovery.
One of its most significant modern pillars is the AdS/CFT correspondence.
This duality relates a theory of gravity in a specific volume (Anti-de Sitter space) to a quantum field theory on its boundary.
This ``holographic principle'' has allowed physicists to use the mathematics of string theory to solve problems in nuclear physics and condensed matter,
even in the absence of direct experimental evidence for strings themselves.

The long-term ambition remains a unified, possibly background-independent framework.
In 1995, Edward Witten proposed that the five consistent superstring theories are different limits of a deeper, 11-dimensional framework known as M-theory.
While M-theory introduced branes and suggested that an additional spatial dimension emerges at strong coupling, it remains an incomplete framework.
We currently lack a non-perturbative master principle—a definitive ``M-theory action''—that would serve as the theory's starting point.

Ultimately, direct experimental tests of string-scale physics remain out of reach.
While the theory's flexibility allows it to accommodate reality, the mechanism required to derive our specific universe from first principles remains undiscovered.
In this sense, string theory remains a monumental mathematical achievement that has yet to prove its physical inevitability.


\section{Emergent and Holographic Approaches}

A more radical class of ideas treats gravity and spacetime as emergent rather
than fundamental. This perspective arose from puzzles at the intersection of
gravity, thermodynamics, and quantum theory. Black holes behave as
thermodynamic objects: they possess an entropy proportional to the area of
their horizon (Bekenstein--Hawking) and emit thermal radiation (Hawking).
These results suggest a deep link between geometry, information, and
statistical mechanics.

The observation that gravitational entropy scales with area rather than
volume led tHooft1993 and Susskind to propose the
idea that the degrees of freedom of a region of spacetime may be fully encoded
on its boundary. In a theory containing gravity, the maximal information
content of a region is bounded by the area of its boundary in Planck units,
not by its volume. This is a surprising result, since in ordinary quantum
field theory the number of degrees of freedom grows with volume.

However, several caveats should be kept in mind. The dualities are established in
specific settings, mainly asymptotically anti-de Sitter (AdS) spacetimes with
known boundary theories. Both AdS and de Sitter (dS) space are idealized
vacuum solutions of Einstein's equations with a constant energy density of
empty space, the cosmological constant, and they differ in its sign. In AdS
the constant is negative: the geometry curves inward, behaves like a box in
which light reaches the edge and returns, and possesses a timelike boundary on
which a dual quantum theory can be defined. In dS the constant is positive:
space expands ever faster, and there is no such boundary, only a cosmological
horizon beyond which signals can never reach an observer.

Neither geometry describes our universe literally, since both are empty; in
AdS/CFT, matter enters as fields living in the bulk and is encoded in the
boundary theory. Our universe is better described by the $\Lambda$CDM model,
which contains matter and radiation together with a cosmological constant (or
something that mimics one). Because the matter density dilutes as space
expands, the universe is expected to approach a de Sitter state in the far
future if dark energy is truly constant, and it also passed through a nearly
de Sitter phase during inflation. Whether dark energy is exactly a cosmological
constant remains an open observational question. The point is therefore not
that we live in de Sitter space, but that the sign of the cosmological
constant in our universe is the opposite of the one for which holography is
best understood. A well-controlled holographic description of de Sitter space
does not yet exist, although proposals such as dS/CFT are being explored.

Applications of holography to quark--gluon plasma or strange metals use models
that share qualitative features with these systems (for example, the
near-universal ratio $\eta/s = 1/4\pi$ of shear viscosity to entropy density in
Einstein-gravity duals, which real quark--gluon plasma approximately matches).
They are not first-principles derivations. Holography is therefore best viewed
as strong evidence that gravity and spacetime \emph{can} be emergent, not as a
demonstration that they \emph{are} in our universe.


\section{The Simulation Hypothesis: Reality as Software}

During recent years so-called Simulation Hypothesis has become increasingly popular. If $3D$ can map to $2D$, why stop there? In software architecture, any $N$-dimensional space is ultimately stored as a one-dimensional bitstring ($N \rightarrow 1$). A sequence of bits has no intrinsic geometry, it is just a set of information.

If the universe is fundamentally informational, it is tempting to conclude that we are merely a program running on some higher-order "hardware." In this view, the strange "quantization" of our world is simply the resolution of the grid, and the speed of light is the clock-speed of the processor.

However, the simulation hypothesis feels like a philosophical "shell game." It merely translates the mystery of existence by one level: if we are a simulation, who simulated the simulators? 

Furthermore, it ignores the staggering Information Cost of reality. Consider the entropy of a single human being. To simulate even a single strand of DNA with perfect fidelity requires tracking billions of quantum interactions. To harvest enough information from a "parent universe" to simulate an entire "child universe" would require a massive thermodynamic overhead.
To simulate a universe would require a computer larger than the universe itself. Even if the simulation used 'lazy loading'—only rendering the parts of reality that are currently being observed, who would pay the electricity bill for that project?

From a developer's perspective: The simulation hypothesis might explain why General Relativity and Quantum Mechanics appear so difficult to unify.
Anyone familiar with software engineering understands the reality of legacy code.
The fundamental incompatibility between the smooth, geometric curves of General Relativity and the discrete, probabilistic jumps of Quantum Mechanics might not be a profound mystery of nature, but simply a case of poor design.
In this light, the universe is a patchwork of modules written by incompetent architects, at different times, with different programming languages.
